\centerline{\bf \Huge Перечислительная комба}

\begin{flushright}
        \scriptsize{Ответ можно оставлять в виду сумму/произведения цэшек}
\end{flushright}

\problem{1}
Будем называть номером последовательность из 6 цифр.

(а) Сколько всего существует различных номеров? А номеров, все цифры которых чётны?

(b) Сколько номеров, в которых любые две соседние цифры различны?

(c) Сколько номеров, все цифры которых различны?

(d) Сколько номеров, все цифры которых имеют одинаковую четность?

(e) Сколько номеров, у которых есть хоть одна нечетная цифра?

(f) Сколько номеров, содержащих цифру 7 и не содержащих цифры 0 ?


\problem{2}
Сколько десятизначных чисел, в которых все цифры различны, и при этом цифры 4 и 5 стоят рядом?

\problem{3}
На конкурсе сладкоежек 7 участников были награждены 20 одинаковыми пирожными и 2 одинаковыми тортами. Каждому досталась хотя бы одна сладость. Сколькими способами могли распределиться награды?

\problem{4}
Все натуральные числа от 1 до 50 выписали подряд слева направо: $123456789101112 \ldots 484950$. Сколько существует способов вычеркнуть почти все цифры полученного числа (кроме четырех), чтобы оставшиеся четыре цифры образовали (без перестановок) число 2026?

\problem{5}
У бельчонка есть 5 орехов, 8 грибов и 11 ягод. Сколькими способами он может выложить все эти предметы в ряд так, чтобы никакие две ягоды не лежали рядом?

\problem{6}
Найдите количество восьмизначных чисел, произведение цифр которых равно 1400. Ответ необходимо представить в виде целого числа.
